\documentclass[a4paper,11pt]{article}
\usepackage[utf8]{inputenc}
\usepackage[T1]{fontenc}
\usepackage[french]{babel}
\usepackage{geometry}
\geometry{left=2cm,right=2cm,top=1.5cm,bottom=1.5cm}
\usepackage{hyperref}
\usepackage{enumitem}

\begin{document}

\begin{flushleft}
    \textbf{AISSAOUI Ismail} \\
    Ingénieur en Intelligence Artificielle \& Data Science \\
    Chlef, Algérie \\
    Tél : +213 660 70 77 96 \\
    Email : \href{mailto:ismail.aissaoui.pro@gmail.com}{ismail.aissaoui.pro@gmail.com} \\
    Portfolio : \href{https://ismail-aissaoui.vercel.app}{ismail-aissaoui.vercel.app}
\end{flushleft}

\begin{flushright}
    À l'attention de \\
    \textbf{M. Baptiste Caramiaux} \\
    \textbf{M. Gonzalo Ramos} \\
    Équipe HCI Sorbonne, ISIR \\
    Sorbonne Université \\
    Fait le \today
\end{flushright}

\vspace{0.5cm}

\noindent \textbf{Objet : Candidature au sujet de thèse : « L'IA générative comme outil pour la pensée : conception et étude des systèmes de soutien cognitif »}

\vspace{0.5cm}

Monsieur Caramiaux, Monsieur Ramos,

\medskip

Actuellement en dernière année de cycle d'ingénieur en Intelligence Artificielle à l'École Supérieure en Informatique (ESI-SBA), c'est avec un vif enthousiasme que je vous soumets ma candidature pour ce projet de thèse. Votre approche de l'IA générative comme « substrat cognitif » visant à augmenter la réflexion humaine plutôt qu’à la remplacer correspond précisément à ma vision du développement des systèmes intelligents : une IA au service de l'humain, préservant son esprit critique et son engagement métacognitif.

Ma formation et mes projets personnels m’ont permis d’acquérir une expertise technique robuste dans les architectures de modèles de langage (Transformers), le RAG (Retrieval-Augmented Generation) et le fine-tuning. À titre d'exemple, mon projet \textit{Geo-RAG} n'est pas une simple interface de requête, mais un outil conçu pour explorer des connaissances spatiales complexes, facilitant la découverte plutôt que l’automatisation. De même, mon travail sur \textit{Recognizini} intègre une boucle de feedback actif (\textit{human-in-the-loop}), un paradigme qui me semble essentiel pour maintenir l'agence humaine face à des systèmes capables d'inférer et d'agir de manière autonome.

Au-delà de la technique, je suis profondément intéressé par l'étude de l'interaction homme-machine (IHM) et l'impact cognitif des outils de travail intellectuel. Je suis convaincu que pour construire des systèmes de soutien cognitif efficaces, il faut allier une ingénierie logicielle rigoureuse (que je maîtrise via PyTorch, FastAPI et React) à une compréhension fine des processus de raisonnement. Mon expérience de stage chez \textbf{Sonatrach} sur les jumeaux numériques m'a appris l'importance de la confiance utilisateur et de l'explicabilité dans des environnements critiques, des enjeux au cœur de votre problématique sur l'érosion de l'engagement métacognitif.

Rejoindre le groupe HCI Sorbonne au sein de l'ISIR serait pour moi l'opportunité idéale de contribuer à un projet à la croisée des sciences cognitives et de l'intelligence artificielle de pointe. Mon profil d'ingénieur, curieux des implications humaines de la technologie et capable de transformer des concepts théoriques en prototypes fonctionnels, me semble être un atout pour mener à bien cette recherche doctorale.

Je vous remercie de l'attention que vous porterez à mon dossier et je me tiens à votre entière disposition pour un entretien afin de vous exposer plus en détail mes motivations.

Je vous prie d'agréer, Monsieur Caramiaux, Monsieur Ramos, l'expression de mes salutations distinguées.

\vspace{0.4cm}

\begin{flushleft}
    \textbf{AISSAOUI Ismail}
\end{flushleft}

\end{document}
